% !TEX root = ../../../main.tex
% 所属：第一部分：单变量分析学 / 第12章：函数项序列与函数项级数
% 节：12.7 可积函数的逼近
% 正文按原文件顺序保存；由 ../chapter12.tex 调用。
\section{可积函数的逼近}
\begin{theorem}
	设 $I=[a,b]$，若 $f$ 在 $[a,b]$ 上黎曼可积，则存在连续函数列
	\[
	\{f_k\}_{k=1}^{\infty}\subset C[a,b],
	\]
	使得
	\[
	\int_a^b |f(x)-f_k(x)|\dd x\to 0,
	\qquad k\to\infty.
	\]
	也就是说，连续函数在可积函数中是稠密的
\end{theorem}

\begin{proof}

		任取 $\varepsilon>0$.由于 $f(x)$ 在 $[a,b]$ 上 Riemann 可积,由 Darboux 可积性准则,存在分划$P:a=x_0<x_1<\cdots<x_m=b,$使得
		\[
		\overline{S}(P)-\underline{S}(P)<\frac{\varepsilon}{2}.
		\]
		
		记
		\[
		\Delta x_j=x_j-x_{j-1},
		\]
		并设
		\[
		M_j=\sup_{x\in[x_{j-1},x_j]}f(x),
		\qquad
		m_j=\inf_{x\in[x_{j-1},x_j]}f(x),
		\]
		其中 $j=1,2,\cdots,m$.于是
		\[
		\overline{S}(P)-\underline{S}(P)
		=
		\sum_{j=1}^{m}(M_j-m_j)\Delta x_j.
		\]
		故
		\[
		\sum_{j=1}^{m}(M_j-m_j)\Delta x_j<\frac{\varepsilon}{2}.
		\]
		
		下面先用阶梯函数逼近 $f(x)$.在每个小区间 $[x_{j-1},x_j]$ 中任取一点$\xi_j\in[x_{j-1},x_j],$
	
		定义阶梯函数
		\[
		s(x)=f(\xi_j),\qquad x\in[x_{j-1},x_j),
		\]
		其中 $j=1,2,\cdots,m$,并令$s(b)=f(\xi_m).$
	
		当 $x\in[x_{j-1},x_j]$ 时,有
		\[
		m_j\leq f(x)\leq M_j,
		\qquad
		m_j\leq f(\xi_j)\leq M_j.
		\]
		于是
		\[
		|f(x)-s(x)|
		=
		|f(x)-f(\xi_j)|
		\leq M_j-m_j.
		\]
		因此
		\[\int_a^b |f(x)-s(x)|\dd x
			=
			\sum_{j=1}^{m}\int_{x_{j-1}}^{x_j}|f(x)-s(x)|\dd x\leq
			\sum_{j=1}^{m}(M_j-m_j)\Delta x_j
			<
			\frac{\varepsilon}{2}.
		\]
		
		这说明 $f(x)$ 可以在积分意义下先由阶梯函数 $s(x)$ 逼近.
		
		下面再用连续函数逼近阶梯函数 $s(x)$.由于 $f(x)$ 在 $[a,b]$ 上 Riemann 可积,所以 $f(x)$ 有界.故存在 $A>0$,使得
		\[
		|f(x)|\leq A,\qquad x\in[a,b].
		\]
		由 $s(x)$ 的定义可知
		\[
		|s(x)|\leq A,\qquad x\in[a,b].
		\]
		
		阶梯函数 $s(x)$ 只可能在有限个分点$x_1,x_2,\cdots,x_{m-1}$处发生跳跃.取充分小的 $\delta>0$,使得区间
		\[
		(x_j-\delta,x_j+\delta),
		\qquad j=1,2,\cdots,m-1,
		\]
		两两不交,并且满足
		\[
		4A(m-1)\delta<\frac{\varepsilon}{2}.
		\]
		
		现在构造连续函数 $g(x)$.在不属于这些小区间的点上,令$g(x)=s(x).$在每个小区间$[x_j-\delta,x_j+\delta]$上,令 $g(x)$ 从 $s(x_j-\delta)$ 线性连接到 $s(x_j+\delta)$.
		
		这样定义出来的 $g(x)$ 在每个小区间内连续,在连接处也连续,且在其余地方等于阶梯函数 $s(x)$.故$g\in C[a,b].$

		又因为 $g(x)$ 只是把 $s(x)$ 的相邻两个阶梯值作线性连接,而 $s(x)$ 的值都在 $[-A,A]$ 中,所以
		\[
		|g(x)|\leq A,\qquad x\in[a,b].
		\]
		于是
		\[
		|g(x)-s(x)|\leq 2A,\qquad x\in[a,b].
		\]
		
		注意到 $g(x)$ 与 $s(x)$ 不相等的地方只包含在
		\[
		\bigcup_{j=1}^{m-1}(x_j-\delta,x_j+\delta)
		\]
		中,这些区间的总长度不超过$2(m-1)\delta.$

		因此
		\[	\int_a^b |g(x)-s(x)|\dd x
			\leq
			2A\cdot 2(m-1)\delta
			=
			4A(m-1)\delta\\
			<
			\frac{\varepsilon}{2}
		\]
		
		由三角不等式,
		\[\int_a^b |f(x)-g(x)|\dd x
			\leq
			\int_a^b |f(x)-s(x)|\dd x
			+
			\int_a^b |s(x)-g(x)|\dd x
			<
			\frac{\varepsilon}{2}+\frac{\varepsilon}{2}
			=
			\varepsilon.
		\]
		
		所以对任意 $\varepsilon>0$,存在连续函数 $g\in C[a,b]$,使得
		\[
		\int_a^b |f(x)-g(x)|\dd x<\varepsilon.
		\]
		
		令 $\varepsilon=\frac1k$,则对每个 $k\in\mathbb{N}_+$,存在 $f_k\in C[a,b]$,使得
		\[
		\int_a^b |f(x)-f_k(x)|\dd x<\frac1k.
		\]
		于是
		\[
		\int_a^b |f(x)-f_k(x)|\dd x\to0,\qquad k\to\infty.
		\]
		
		综上,$C[a,b]$在 $\mathcal{R}[a,b]$ 中稠密.

\end{proof}
\begin{corollary}\label{C在R中一致有界稠密}
	设 $f(x)$ 在闭区间 $[a,b]$ 上 Riemann 可积,且
	\[
	|f(x)|\leq B,\qquad x\in[a,b].
	\]
	则存在连续函数列
	\[
	\{f_k\}_{k=1}^{\infty}\subset C[a,b],
	\]
	使得
	\[
	\sup_{x\in[a,b]}|f_k(x)|\leq B,\qquad k=1,2,\cdots,
	\]
	并且
	\[
	\int_a^b |f(x)-f_k(x)|\dd x\to0,\qquad k\to\infty.
	\]
\end{corollary}

\begin{proof}
	由上一定理可知,存在连续函数列
	\[
	\{h_k\}_{k=1}^{\infty}\subset C[a,b],
	\]
	使得
	\[
	\int_a^b |f(x)-h_k(x)|\dd x\to0,\qquad k\to\infty.
	\]
	
	但是此时 $h_k(x)$ 未必满足
	\[
	\sup_{x\in[a,b]}|h_k(x)|\leq B.
	\]
	因此需要对 $h_k(x)$ 作截断.
	
	定义函数
	\[
	\varphi(t)=
	\begin{cases}
		-B, & t<-B,\\
		t, & -B\leq t\leq B,\\
		B, & t>B.
	\end{cases}
	\]
	显然 $\varphi(t)$ 是 $\mathbb{R}$ 上的连续函数,并且$|\varphi(t)|\leq B,\qquad t\in\mathbb{R}$
	
	令
	\[
	f_k(x)=\varphi(h_k(x)),\qquad x\in[a,b].
	\]
	由于 $h_k(x)$ 连续,$\varphi(t)$ 连续,所以$f_k\in C[a,b].$
	
	又由 $|\varphi(t)|\leq B$ 可知
	\[
	|f_k(x)|\leq B,\qquad x\in[a,b].
	\]
	于是
	\[
	\sup_{x\in[a,b]}|f_k(x)|\leq B.
	\]
	
	下面证明
	\[
	\int_a^b |f(x)-f_k(x)|\dd x\to0.
	\]
	
	由于
	\[
	|f(x)|\leq B,\qquad x\in[a,b],
	\]
	所以$f(x)\in[-B,B].$而 $\varphi(t)$ 是把实数 $t$ 截断到区间 $[-B,B]$ 上的函数,所以对任意 $t\in\mathbb{R}$ 和任意 $u\in[-B,B]$,都有
	\[
	|\varphi(t)-u|\leq |t-u|.
	\]
	取
	\[
	t=h_k(x),\qquad u=f(x),
	\]
	便得到
	\[
	|\varphi(h_k(x))-f(x)|
	\leq
	|h_k(x)-f(x)|.
	\]
	也就是
	\[
	|f_k(x)-f(x)|
	\leq
	|h_k(x)-f(x)|.
	\]
	
	两边在 $[a,b]$ 上积分,得
	\[
	\int_a^b |f_k(x)-f(x)|\dd x
	\leq
	\int_a^b |h_k(x)-f(x)|\dd x.
	\]
	由
	\[
	\int_a^b |h_k(x)-f(x)|\dd x\to0
	\]
	可知
	\[
	\int_a^b |f_k(x)-f(x)|\dd x\to0.
	\]
	
	综上,存在连续函数列 $\{f_k\}_{k=1}^{\infty}\subset C[a,b]$,满足
	\[
	\sup_{x\in[a,b]}|f_k(x)|\leq B,
	\]
	并且
	\[
	\int_a^b |f(x)-f_k(x)|\dd x\to0.
	\]
\end{proof}